\section{Discussion}
\label{sec:discussion}

\paragraph{What the theory establishes.}
Exact information reduction has a precise hierarchy.  At the metric level,
submetry is necessary and sufficient for universal radial-visible reduction.
Coherent sections select reusable information sheets.  Split-Hadamard
geometry turns those sheets into global horizontal submanifolds.  SPD
compression realizes the entire hierarchy in closed form and adds a rank
quotient, a closed-form completion-pair metric, and computable structured
certificates in finite-dimensional LP/SDP feasibility form.

\paragraph{Visibility and reference.}
A visible target does not determine an ambient geometry.  The invisible fiber
has infinite diameter, and even its radius-bounded minimax problem has risk
equal to the radius.  The canonical completion is therefore meaningful only
relative to a declared \(P_0\).  Its role is precise: preserve the reference
on directions not forced to change by the visible target while realizing the
visible target with minimum AIRM deformation.

\paragraph{Moving information.}
The mismatch weight \(\Xiop(R)=R+R^{-1}-2I\) explains when subspace motion is
observable in completed geometry.  Rotating a mode whose visible geometry
equals the reference changes no completed matrix.  As a mode separates from
the reference, the same rotation acquires positive cost.  Continuous rank
changes can therefore occur only through reference-valued modes.  These are
kinematic statements; choosing a trajectory requires an additional modeling,
optimization, or control principle.  The bilinear kernel also separates two
second-order objects.  The reduced minimum-distance value is independent of
the normalized subspace coordinate at fixed \(R\), while the minimizing
completed matrix generally moves with that coordinate.  Its motion is
measured by the positive-semidefinite Gram form
\eqref{eq:completion-pair-kernel}.  Reduced-value curvature and pullback
metric geometry therefore answer different questions.

\paragraph{Structured optimizers.}
Diagonal and block geometry are restrictions of the same visible map
\(P\mapsto A^\top PA\).  Their relative-interior tests distinguish three
scientifically different outcomes: strict representability, boundary
representability requiring degeneracy, and infeasibility even after closure.
A single residual value cannot make these distinctions; the conic certificates
can.

\paragraph{Statistical interpretation.}
Known-subspace recovery is intrinsically \(m\)-dimensional.  Estimating the
subspace introduces an ambient second-moment problem and an eigengap.  The exact
line element then converts principal-angle error into completed AIRM error
with a geometry-dependent weight.  The resulting bound separates statistical
uncertainty in \(R\) from uncertainty in the information channel itself.  The
self-contained rates use bounded i.i.d. observations and net--Hoeffding
arguments; they establish consistency in the stated regime and are not
rate-optimal claims.

\paragraph{Boundaries.}
The abstract theorem characterizes exactness but does not make an arbitrary
observation map a submetry.  Nonlinear, infinite-dimensional, multi-view, or
non-AIRM models require separate verification.  The SPD pair metric is
specific to full-column-rank linear compression and is analytic across
repeated eigenvalues.  The structured certificate
theorems cover diagonal and block families; Kronecker, low-rank, and tensor
families have different image geometry.  Finally, the paper assumes a
positive visible target.  Indefinite curvature must be damped, projected, or
otherwise converted before this theory applies.

\paragraph{Downstream use.}
The paper supplies a geometric foundation for later questions about metric
evolution, constrained metric paths, and intervention-based optimizer audits.
Those questions may use the visible state \((A,R)\), the canonical completion
\(\Psi_{P_0,A}(R)\), and the completion-pair metric derived here.  Their
dynamical or causal assumptions are intentionally outside the present paper.
